% Vietnamese translation for OI Public Library
% Source-File: IOI中国国家候选队论文/2003/雷环中/雷环中.doc
\documentclass[11pt]{article}
\usepackage{oipl-vn}

\title{Bàn về bài toán dạng nộp kết quả}
\author{Lei Huanzhong}
\date{2003}

\begin{document}
\maketitle

\origtitle{On output-only problems}
\sourcepath{IOI China candidate team papers / 2003 / Lei Huanzhong / Lei Huanzhong.doc}

\begin{abstract}
Bài toán dạng nộp kết quả có lịch sử còn rất ngắn, nhưng nhờ mô hình mới,
cách nhìn mới về hiệu quả, thời gian và bộ nhớ, cùng các kỹ thuật giải độc
đáo, nó nhanh chóng trở thành một dạng bài quan trọng trong Olympic Tin học.
Bài viết thảo luận hai kỹ thuật chính của dạng bài này: phân tích dữ liệu và
ngẫu nhiên hóa. Nguyên tắc cốt lõi có thể tóm lại bằng một câu: trong thời
gian ngắn nhất, hãy lấy được kết quả đúng; điều quan trọng không phải là quá
trình.
\end{abstract}

\noindent\textbf{Từ khóa:} nộp kết quả, phân tích dữ liệu, ngẫu nhiên hóa,
chiến lược.

\section{Dẫn nhập}

Bài toán dạng nộp kết quả là một trong những dạng bài trẻ nhất của Olympic
Tin học. Việc IOI liên tiếp xuất hiện dạng bài này trong hai năm cho thấy nó
đã bước vào các kỳ thi chính thống. Lý do quan trọng là dạng bài này linh
hoạt, mới mẻ, khuyến khích nhiều con đường khác nhau đi tới cùng một kết quả,
và tạo không gian cho những phương pháp khác thường.

Nhìn lại các bài như \textit{Crack the Code}, \textit{Double Crypt},
\textit{Birthday Party}, \textit{Tetris} và \textit{Xor}, ta thấy mỗi bài đều
có ý tưởng thú vị, và phương pháp giải cũng rất đa dạng. Nhưng muốn làm đúng
và nhanh dạng bài này, ta cần suy nghĩ sâu hơn về chính đặc điểm của nó.

Sự xuất hiện của bài nộp kết quả làm thay đổi nhiều thứ.

\paragraph{Thay đổi nội dung khảo sát.}
Khi Olympic Tin học đã phát triển tương đối chín, nếu chỉ khảo sát thuật toán
và cấu trúc dữ liệu thì chưa đủ. Bài nộp kết quả còn khảo sát năng lực tư duy,
khả năng sáng tạo, chiến lược tổng thể và chiến lược thời gian. Những yếu tố
này cũng tồn tại trong các bài truyền thống, nhưng ở dạng nộp kết quả chúng
được thể hiện trực tiếp hơn nhiều.

\paragraph{Thay đổi độ phân hóa.}
Với bài truyền thống, ta thường so sánh thuật toán bằng độ phức tạp thời gian
và bộ nhớ. Nhưng ở bài nộp kết quả, điều quan trọng hơn là thời gian để tạo
ra các tệp đáp án đúng. Nếu hai người đều tạo đáp án đúng, người hoàn thành
trong 20 phút rõ ràng thắng người mất 2 giờ, bất kể họ dùng phương pháp nào.

\paragraph{Thay đổi bản chất đề bài.}
Một bài như \textit{Birthday Party} nếu được ra dưới dạng nộp chương trình
thì cách nghĩ sẽ khác hẳn so với dạng nộp kết quả. Khi tác giả đề bài chọn
dạng nộp kết quả, họ không muốn thí sinh chỉ nhìn nó như một bài chuẩn, mà
muốn thí sinh tận dụng chính dữ liệu đầu vào đã cho, mạnh dạn dùng mọi cách
để thu được kết quả đúng nhanh nhất.

\paragraph{Thay đổi cách chấm.}
Chấm bài lập trình vốn đã là hộp đen; với bài nộp kết quả, ``hộp đen'' càng
nhấn mạnh vào kết quả cuối cùng. Chương trình hay quá trình sinh ra đáp án
không còn quan trọng bằng việc tệp nộp có đúng hay không.

Vì vậy, nguyên tắc giải bài có thể viết ngắn gọn:

\begin{center}
\textit{Luôn nhắc mình rằng thứ cần có là kết quả đúng trong thời gian ngắn
nhất, không phải một quá trình hoàn hảo.}
\end{center}

\section{Phân tích dữ liệu}

Phân tích dữ liệu gần như sinh ra cùng với bài nộp kết quả. Nó đòi hỏi khả
năng quan sát, phỏng đoán, thực nghiệm và thao tác thủ công cao hơn nhiều so
với các phương pháp thông thường. Phần này dùng hai ví dụ:
\textit{Crack the Code} và \textit{Birthday Party}.

\subsection{\textit{Crack the Code} (Baltic OI 2001)}

\paragraph{Tóm tắt đề.}
Có năm cặp tệp: mỗi tệp \texttt{cra*.txt} là một văn bản gốc tiếng Anh, và
tệp \texttt{cra*.in} tương ứng được tạo từ một tệp chưa biết
\texttt{cra*.out}. Với mỗi tệp, có 10 số kiểu byte chưa biết. Khi xử lý ký tự
thứ \(i\), nếu ký tự không phải chữ cái thì giữ nguyên; nếu là chữ thường thì
đổi thành chữ hoa; nếu là chữ hoa thì dịch vòng trong bảng chữ cái theo một
số bước phụ thuộc vào \(i\bmod 10\). Nhiệm vụ là từ \texttt{cra*.in} và
\texttt{cra*.txt} khôi phục \texttt{cra*.out}.

Bài này không có một thuật toán thật sự hiệu quả cho mọi dữ liệu, vì bản chất
là một dạng mã hóa khá hoàn chỉnh. Đó chính là lý do nó phù hợp với hình thức
nộp kết quả: chỉ cần phá được các dữ liệu đã cho.

\paragraph{Dùng dữ liệu để phá mã.}
Điểm đầu tiên cần chú ý là chỉ có 5 bộ dữ liệu, và mỗi văn bản đều là tiếng
Anh. Trong tiếng Anh, nhiều từ tần suất cao như \textit{a}, \textit{I},
\textit{the}, \textit{he}, \textit{are}, \textit{of}, \textit{in},
\textit{after} rất ngắn; từ càng ngắn thì khả năng càng ít.

Mục tiêu chính là tìm 10 số dịch chuyển. Nếu xác định được một vài cặp từ
tương ứng giữa bản rõ và bản mã, các dịch chuyển sẽ lộ ra ngay.

Trong dữ liệu thứ nhất, đoạn mã có cụm \texttt{EQMDY 1943}. Dựa vào văn bản
gốc và ngữ cảnh tiểu sử, gần như chắc chắn \texttt{EQMDY} là
\texttt{AFTER}. Chỉ từ một từ này đã suy ra được 5 trong 10 số dịch chuyển.
Sau khi thử khôi phục bằng 5 số đó, nhiều cụm từ tiếp theo hiện ra gần đúng,
chẳng hạn \texttt{THG NXHPARD UNKBBHMITY} rõ ràng là
\texttt{THE HARVARD UNIVERSITY}, và \texttt{UNIVETYFJS OF CALKLLHHIA} là
\texttt{UNIVERSITY OF CALIFORNIA}. Dùng thêm một hai cụm dài như vậy là đủ
khôi phục toàn bộ văn bản.

Trong dữ liệu thứ năm, cách thủ công còn nhanh hơn. Văn bản gốc nói về một
không gian vectơ tô pô; dòng mã bắt đầu bằng một từ 10 chữ cái kèm dấu hai
chấm. Rất tự nhiên đoán đó là \texttt{DEFINITION}. Thử đoán này thì toàn bộ
đoạn trở nên rõ:

\begin{verbatim}
DEFINITION: THE DUAL SPACE OF A TOPOLOGICAL VECTOR SPACE X
IS THE VECTOR SPACE X* WHOSE ELEMENTS ARE THE CONTINUOUS
LINEAR FUNCTIONALS ON X.
\end{verbatim}

Ở đây, trực giác ngôn ngữ và tri thức nền của con người đóng vai trò quyết
định. Nếu các bài truyền thống thường khảo sát thí sinh thông qua chương
trình, thì dạng bài này gần như đối thoại trực tiếp với thí sinh.

\paragraph{Phương pháp thống kê tần suất.}
Ta cũng có thể dùng chương trình. Vì văn bản tiếng Anh chỉ có 26 chữ cái, tần
suất xuất hiện của mỗi chữ không đều. Gọi \(P\) là phân bố tần suất chữ cái
trong văn bản gốc. Với mỗi lớp vị trí \(i,i+10,i+20,\ldots\) trong văn bản mã,
tính phân bố \(H\), rồi thử mọi dịch chuyển \(0,\ldots,25\) và chọn dịch
chuyển làm \(H\) giống \(P\) nhất.

Bài gốc so sánh vài hàm đánh giá mức giống nhau. Kết quả cho thấy khi mẫu đủ
dài, phương pháp thống kê khá tốt; khi văn bản quá ngắn, tần suất mất ý nghĩa
và phương pháp thất bại. Khi đó dữ liệu phân tích thủ công lại có ưu thế.

Điều rút ra là: với bài nộp kết quả, không có phương pháp nào luôn tốt nhất.
Phân tích dữ liệu, chương trình phụ trợ ngắn, thống kê tần suất và kiểm tra
thủ công thường phải phối hợp với nhau.

\subsection{\textit{Birthday Party} (CEOI 2002)}

\paragraph{Tóm tắt đề.}
Có \(N\) người và \(M\) yêu cầu lô-gic. Một yêu cầu có thể là tên một người,
phủ định của tên đó, phép \(\land\), phép \(\lor\), hoặc phép kéo theo
\(\Rightarrow\) giữa các yêu cầu. Với 10 tệp \texttt{party*.in}, cần chọn một
tập người sao cho mọi yêu cầu đều đúng. Đề bảo đảm có ít nhất một nghiệm và
chỉ cần nộp một nghiệm. Kích thước có thể tới
\[
  N\le 25000,\qquad M\le 129998.
\]

Nhìn hình thức, đây là SAT, vốn là bài toán NP-đầy đủ. Quy mô dữ liệu cho
thấy nếu giải được thì mấu chốt chắc chắn nằm ở dữ liệu cụ thể, không nằm ở
một thuật toán tổng quát cho SAT.

\paragraph{Phân tích từng dữ liệu.}
Dữ liệu đầu tiên rất nhỏ, có thể giải bằng tay hoặc bằng chương trình đơn
giản. Với dữ liệu thứ hai và thứ ba, biểu thức hầu hết có dạng mệnh đề OR
ngắn; tên biến lại có quy luật rất rõ: \texttt{a,b,c,...,j,ba,bb,...} tương
ứng tự nhiên với các số liên tiếp. Chỉ cần một chương trình quét dữ liệu là
có thể phát hiện quy luật này và xử lý thuận lợi hơn.

Dữ liệu thứ tư có vài dòng biểu thức phức tạp hơn, nhưng số dòng đặc biệt rất
ít. Một chiến lược hợp lý là dùng chương trình xử lý phần lớn mệnh đề đơn
giản, rồi kiểm tra thủ công sáu yêu cầu đặc biệt. Viết một bộ phân tích cú
pháp đầy đủ chỉ để phục vụ sáu dòng như vậy thường không đáng.

Dữ liệu thứ năm đến thứ bảy hầu như toàn là dạng
\[
  A\Rightarrow B,
\]
trong đó \(A,B\) là biến hoặc phủ định biến. Đây là 2-SAT dưới dạng đồ thị
hàm ý. Dữ liệu thứ tám và thứ chín hầu như toàn là dạng
\[
  A\lor B,
\]
cũng có thể chuyển về 2-SAT. Dữ liệu thứ chín còn có vài dòng đặc biệt, nhưng
số lượng rất ít và có thể xử lý riêng. Dữ liệu thứ mười ban đầu có vẻ rất
đáng sợ vì biểu thức không đều, nhưng quan sát phần cuối sẽ thấy nhiều ràng
buộc dạng
\[
  x\Rightarrow \neg x,\qquad \neg y\Rightarrow y,
\]
ép trực tiếp giá trị của rất nhiều biến. Phần còn lại vì thế nhỏ hơn nhiều.

Bài học ở đây là: dữ liệu không chỉ là đầu vào, mà còn là gợi ý chiến lược.
Nếu không quan sát kỹ, ta dễ bỏ qua vài dòng đặc biệt hoặc vài quy luật đặt
tên biến; nhưng chính những chi tiết ấy quyết định cách giải nhanh nhất.

\paragraph{So sánh với bài thường.}
Nếu \textit{Birthday Party} là bài nộp chương trình, ta phải xây dựng một bộ
xử lý biểu thức lô-gic khá tổng quát. Nhưng với bài nộp kết quả, có thể kết
hợp nhiều thủ đoạn: giải tay dữ liệu nhỏ, dùng 2-SAT cho dữ liệu có cấu trúc,
viết chương trình quét để phát hiện dòng đặc biệt, và kiểm tra thủ công vài
biểu thức ngoại lệ.

\section{Ngẫu nhiên hóa}

Ngẫu nhiên hóa cũng rất phù hợp với bài nộp kết quả, vì ta có thể chạy nhiều
lần, chọn kết quả tốt nhất, và không cần chứng minh quá trình luôn ổn định
như trong bài nộp chương trình. Hai ví dụ tiêu biểu là \textit{Tetris} và
\textit{Xor}.

\subsection{\textit{Tetris}}

Trong bài \textit{Tetris}, mỗi dữ liệu cho một trạng thái ban đầu và cần tạo
một chuỗi thao tác đặt khối để làm phẳng bàn chơi. Một số dữ liệu nhỏ có thể
tính tay, một số có quy luật rõ chỉ cần chương trình ngắn, còn các dữ liệu
lớn không thể làm thủ công.

Một thuật toán khả thi là dùng vài dạng khối ``dễ chịu'' để trước hết đưa bàn
về trạng thái thấp, sau đó dùng ngẫu nhiên hóa để lấp các chỗ trống từ trái
sang phải. Mỗi lần gặp một vị trí cần xử lý, chọn ngẫu nhiên một khối hợp lệ
để lấp. Thuật toán này thô nhưng cài rất nhanh.

Sau đó có thể thêm một số kinh nghiệm thủ công: khi bàn đã khá phẳng, ưu tiên
các khối làm bề mặt tiếp tục phẳng; tránh các khối tạo ra độ nhô đột ngột; ưu
tiên vài khối đặc biệt nếu chúng làm giảm chiều cao tổng thể. Những quy tắc
này không phải chứng minh chặt, nhưng trong bài nộp kết quả chỉ cần chúng tạo
đáp án đúng cho dữ liệu đã cho.

Bài gốc còn nêu hai phương pháp ổn định hơn:

\begin{itemize}
  \item Mỗi bước chọn thao tác làm bàn phẳng nhất theo một hàm đánh giá như
  tổng độ lệch giữa chiều cao các cột.
  \item Chia bàn thành nhóm 4 cột, xử lý từng nhóm bằng cấu trúc xây dựng,
  rồi ghép các nhóm lại.
\end{itemize}

Phương pháp xây dựng có chất lượng tốt và chạy nhanh, nhưng cài phức tạp hơn.
Ngược lại, ngẫu nhiên hóa dễ viết, đủ nhanh để thử nhiều lần, và phù hợp với
mục tiêu lấy kết quả trong thời gian ngắn.

\subsection{\textit{Xor} (IOI 2002)}

Bài \textit{Xor} cho một bảng \(N\times N\) tô đen trắng. Mỗi thao tác chọn
một hình chữ nhật và đảo màu mọi ô trong đó. Cần dùng càng ít thao tác càng
tốt để biến toàn bộ bảng thành trắng. Mỗi dữ liệu được tính điểm theo số thao
tác.

Kinh nghiệm cho thấy nên nhìn vào các ``góc'' của vùng đen. Có thể chia góc
thành ba loại; góc loại một và hai chắc chắn là góc của một hình chữ nhật
trong một lời giải nào đó, còn góc loại ba thì không nhất thiết. Thuật toán
tham lam cơ bản:

\begin{enumerate}
  \item Tạo tập \(A\) gồm mọi góc loại một và hai.
  \item Nếu trong \(A\) có bốn điểm là bốn đỉnh của một hình chữ nhật, xem đó
  là một thao tác và xóa bốn điểm khỏi \(A\).
  \item Nếu không, chọn ba điểm tạo thành ba đỉnh của một hình chữ nhật, thêm
  đỉnh thứ tư còn thiếu và thực hiện thao tác tương ứng; cập nhật \(A\).
  \item Lặp đến khi \(A\) rỗng.
\end{enumerate}

Phần xóa bốn điểm không phụ thuộc nhiều vào thứ tự, nhưng phần chọn ba điểm
có ảnh hưởng rõ tới chất lượng lời giải. Vì \(N\le 2000\), thử một chiến lược
chọn phức tạp có thể tốn thời gian; do đó có thể chọn ngẫu nhiên nhiều lần và
lấy lời giải tốt nhất. Thực nghiệm trong bài gốc cho thấy sau nhiều lần chạy,
kết quả rất gần lời giải tốt nhất đã biết, một số dữ liệu còn đạt mức tối ưu
đã biết.

Đây là một ưu thế đặc thù của nộp kết quả: ta không cần một thuật toán có
hiệu năng ổn định trên mọi lần chạy; ta chỉ cần có một lần chạy sinh ra tệp
đáp án tốt.

\section{Chiến lược cho bài nộp kết quả}

Dạng bài này không chỉ khảo sát thuật toán, mà còn khảo sát chiến lược thi.

\subsection{Cách nhìn mới về hiệu quả}

Trong bài truyền thống, ta thường hình dung quá trình làm bài như sau: ban
đầu mất thời gian thiết kế thuật toán và cấu trúc dữ liệu, sau đó chương trình
mới dần đạt hiệu quả cao. Có những lúc ta tiếp tục tối ưu thuật toán dù kết
quả hiện tại đã đủ để qua bài; về chiến lược thi, đó có thể là lãng phí.

Với bài nộp kết quả, đường cong thường khác: ban đầu có thể lấy điểm rất
nhanh bằng phân tích dữ liệu, giải tay hoặc chương trình phụ trợ nhỏ; về sau
mỗi điểm tăng thêm càng tốn công hơn. Vì dữ liệu đã cho trước, ta có thể ước
lượng khá chính xác mình còn lấy được bao nhiêu điểm trong bao lâu. Do đó
cần quyết định khi nào nên dừng, khi nào nên bỏ một vài dữ liệu khó để dành
thời gian cho bài khác, và khi nào nên viết một chương trình tổng quát hơn.

\subsection{Cách nhìn mới về thời gian và bộ nhớ}

Bài nộp kết quả phá vỡ nhiều thói quen của bài truyền thống:

\begin{itemize}
  \item bộ nhớ dùng khi sinh đáp án thường rộng rãi hơn;
  \item thời gian chạy có thể dài hơn nhiều, thậm chí vài giờ nếu cuộc thi cho
  phép;
  \item có thể chạy nhiều tiến trình hoặc nhiều thử nghiệm song song;
  \item dữ liệu đầu vào cố định, nên có thể khai thác mọi đặc điểm riêng của
  nó.
\end{itemize}

Điều này không có nghĩa bỏ qua thuật toán đẹp, mà là phải kết hợp thực lực
với mẹo phù hợp. Mục tiêu là tệp kết quả đúng, không phải một chương trình
đẹp cho mọi dữ liệu tưởng tượng.

\subsection{\textit{Double Crypt} (IOI 2001)}

Bài \textit{Double Crypt} dùng mã hóa AES hai lần:
\[
  c=E(E(p,k_1),k_2),
\]
với \(k_1,k_2\) nằm trong một tập khóa giới hạn bởi tham số \(s\). Cần tìm
\(k_1,k_2\) từ \(p,c,s\).

Thuật toán chuẩn là meet-in-the-middle: tính mọi \(E(p,k_1)\), lưu vào bảng
băm; đồng thời tính mọi \(D(c,k_2)\) và tìm giao. Nếu bỏ qua va chạm, thời
gian và bộ nhớ đều khoảng \(O(2^s)\). Thuật toán này rõ ràng tốt nhưng cần
cài đặt cẩn thận.

Trong bối cảnh nộp kết quả, có thể tạm dùng cách liệt kê đơn giản hơn cho
một số dữ liệu, tận dụng các đặc điểm như dữ liệu thứ ba và thứ chín có cùng
\(p,c\), hoặc một số dữ liệu có thể chạy chung một lần. Dù chương trình chạy
lâu hơn, ta có thể để nó chạy nền trong khi làm bài khác. Cách này không hoàn
hảo, nhưng có thể nhanh chóng lấy phần lớn điểm, phù hợp với chiến lược tổng
thể.

\subsection{Phối hợp tính tay và chương trình}

Nhiều bài nộp kết quả hay ở chỗ buộc ta phối hợp tính tay với chương trình.
\textit{Crack the Code} và \textit{Birthday Party} đều như vậy. Nếu quá lệ
thuộc vào tính tay, các dữ liệu lớn sẽ không làm nổi; nếu quá lệ thuộc vào
chương trình tổng quát, ta mất thời gian cài những thứ không cần thiết cho
vài dòng ngoại lệ.

Người giải giỏi là người biết dùng cả bộ não lẫn máy tính: dùng chương trình
ngắn để quét, thống kê, kiểm tra; dùng mắt người để nhận ra quy luật, đoán từ
khóa, xử lý ngoại lệ; rồi lại dùng chương trình để xác minh và xuất tệp.

\section{Kết luận}

Bài toán dạng nộp kết quả buộc ta rời khỏi cách nhìn truyền thống. Phải phân
tích từng dữ liệu cụ thể, thiết kế phương pháp từ góc nhìn toàn cục, và phối
hợp linh hoạt giữa thuật toán, quan sát, thử nghiệm, ngẫu nhiên hóa, tính tay
và chương trình phụ trợ.

Cùng một bài toán, nếu là nộp chương trình thì có thể cần một thuật toán ổn
định, tổng quát; nếu là nộp kết quả thì phương pháp tốt nhất có thể là một
tổ hợp các thủ đoạn nhỏ. Tiêu chuẩn đánh giá đã thay đổi, nên chiến lược cũng
phải thay đổi.

Tóm lại: hãy luôn nhắc mình rằng thứ cần có là kết quả đúng trong thời gian
ngắn nhất, không phải quá trình.

\begin{thebibliography}{9}
\bibitem{boi2001}
Marcin Kubica.
\textit{The 7th Baltic Olympiad in Informatics BOI 2001 Tasks and Solutions}.

\bibitem{ceoi2002}
\textit{Summary of the 9th Central European Olympiad in Informatics}.

\bibitem{noi2002}
\textit{Problems of the 19th National Olympiad in Informatics}, 2002.

\bibitem{ioi2002}
\textit{IOI 2002 Problems}.

\bibitem{ioi2001}
Jyrki Nummenmaa, Erkki Makinen, Isto Aho.
\textit{IOI'01 Competition}, second edition.
\end{thebibliography}

\end{document}
